\chapter{ADHD Screening}

\section*{What Is This Test?}
ADHD screening is a structured assessment for symptoms of attention-deficit/hyperactivity disorder. It is not a blood test or a single diagnostic examination. Screening identifies people who need a more complete clinical evaluation.

\section*{Alternative Names}
\begin{itemize}
  \item Attention-Deficit/Hyperactivity Disorder Screening
  \item ADHD Questionnaire
  \item ADHD Rating Scale
  \item Attention and Hyperactivity Assessment
\end{itemize}

\section*{Principle}
Clinicians use validated rating scales together with developmental history, interviews, observation, and reports from more than one setting. Diagnosis requires a persistent pattern of symptoms that causes impairment, began during development, and is not better explained by another condition.

\section*{Why This Test Is Ordered}
Screening is considered for concerns about inattention, impulsivity, hyperactivity, school or work performance, relationships, organization, or unsafe behavior. Sleep problems, anxiety, depression, learning disorders, substance use, and medical conditions may produce similar symptoms.

\section*{Primary vs Secondary Test}
A rating scale is a screening aid, not a standalone diagnostic test. The primary evaluation is a comprehensive assessment by a qualified clinician using DSM criteria and collateral information when available.

\section*{How This Test Is Performed}
The patient and, when appropriate, parents, teachers, or partners complete questions about symptoms and functional impairment. The clinician reviews records, developmental history, mental health, sleep, substance use, and medical factors.

\section*{What Preparation Is Needed}
No physical preparation is required. Bring school or work reports, prior evaluations, medication information, and examples of symptoms across settings.

\section*{Understanding Results}
A positive screen means that further evaluation is warranted; it does not confirm ADHD. A negative screen lowers the likelihood but may not exclude ADHD when symptoms are subtle or the informant has limited information.

\section*{Follow-up Tests}
Follow-up may include psychoeducational testing, hearing or vision assessment, sleep evaluation, mental-health assessment, or targeted medical tests when another explanation is suspected. There is no routine laboratory test that confirms ADHD.

\section*{References}
\begin{enumerate}
  \item MedlinePlus. ADHD Screening.
  \item American Academy of Pediatrics. Clinical Practice Guideline for the Diagnosis, Evaluation, and Treatment of ADHD.
  \item American Psychiatric Association. Diagnostic and Statistical Manual of Mental Disorders.
\end{enumerate}

\clearpage
\section*{Understanding Results and Follow-up}
The laboratory report should identify the specimen, method, units, reference interval or decision limit, and whether the result is positive, negative, abnormal, or inconclusive. Reference intervals vary with age, sex, pregnancy, specimen type, assay, and laboratory; a result outside a range is not automatically a diagnosis, and a result within a range does not always exclude disease. Discuss unexpected results with the ordering clinician, who can decide whether repeat or confirmatory testing, treatment, or referral is appropriate. Do not change medicines or delay urgent care based on an isolated result.


\section*{Regulatory Status and Authorization Date}
This chapter describes a test category, not a specific commercial product. FDA, CDSCO, EU IVDR, and MHRA approval or registration dates therefore cannot be assigned without the assay manufacturer, product name, instrument, and intended use. For laboratory-developed tests, an individual agency approval date may not exist. See the Preface for the interpretation of regulatory status.

\clearpage
